\documentclass[10pt,a4paper]{amsart}
\usepackage[margin=19mm]{geometry}
\usepackage{amsmath,amssymb,mathtools}
\usepackage[scr=rsfs,cal=euler]{mathalfa}
\usepackage{xfrac}
\usepackage[unicode,hidelinks,bookmarks=false]{hyperref}
\newcommand{\ie}{i.e.\ }
\newcommand{\cf}{cf.\ }
\newcommand{\resp}{respectively\ }
\newcommand{\Subsection}{\subsection}
\providecommand{\nobreakdash}{\nobreak\hbox{-}}
\setlength{\emergencystretch}{2em}
\multlinegap0pt
\usepackage{amssymb}
\usepackage{mathtools}
\usepackage{xcolor}
\usepackage{tikz}
\usepackage[shortlabels]{enumitem}
\newtheorem{theorem}{Theorem}
\newtheorem{lemma}{Lemma}
\title{Local large deviations for linear-region growth in random piecewise-linear networks}
\author{Recep \"Ozkan, Christian Hirsch}
\date{Source: arXiv:2607.07014v1 (CC BY 4.0)}
\begin{document}
\maketitle

\noindent\emph{Source and licence.} Local large deviations for linear-region growth in random piecewise-linear networks. Version arXiv:2607.07014v1 (CC BY 4.0), distributed by arXiv under the Creative Commons Attribution 4.0 International licence, http://creativecommons.org/licenses/by/4.0/, as recorded in the arXiv licence metadata for this exact version and shown on the abstract page https://arxiv.org/abs/2607.07014v1. Full licence notice: \texttt{licenses/arxiv-2607.07014-CC-BY-4.0.txt}.\par\smallskip

\noindent\textit{Editorial scope.} Counted unit: the article's theorem thm:main-bridge-lower (source0.tex lines 647--650). The article's complete original proof of it is delivered with it (source0.tex lines 1077--1131, one \texttt{proof} environment of 55 lines, which the article places away from the statement). No mathematics is written, repaired or re-derived: the statement and the proof are the article's own lines, in the article's own notation, and no retained line is substituted or re-flowed.  Retained uncounted as the source-local context the proof needs: lines 994--1002.  Nothing was omitted from any retained span, so the package carries no omission record.  No retained line contains a citation, so the wrapper carries no bibliography and no bibliography entry.  No retained line contains a figure environment, an image-inclusion command or drawing code, so no image is delivered and the package declares no asset dependency.  Editorial formatting only, in addition: an AMS \texttt{amsart} preamble that restates the article's own declarations and macros used by the retained lines, namely the environments \texttt{theorem}, \texttt{lemma} on separate counters, together with the standard packages those lines need.  The retained environments print in this document as Theorem 1, Lemma 1 in order of appearance, whereas the article prints the same items with the numbering of its own full document.  The complete native source of the article as distributed in the arXiv e-print for this exact version is bundled unchanged as \texttt{sources/204676/source0.tex}.
\begin{lemma}\label{thm:bridge-lemma}
For all words \(u,v\) in the coarse alphabet, we have
\(
\Gamma\bigl(v b^{(m)}u\bigr)\ge \Gamma(v)\Gamma(u),
	\)
with	\(
 b^{(m)}=(m).
\)
\end{lemma}

\begin{theorem}[Bridged lower bound for the upper tail]\label{thm:main-bridge-lower}
Fix \(m,r\ge1\), and assume \(p_b>0\). Let \(n_k:=1+kr+(k-1)\). Then, for every \(x\in\mathbb R\), one has 
\[\liminf_{k\to\infty} n_k^{-1}\log\mathbb P(n_k^{-1}\log N_{n_k}>x)\ge -J_{r,m}^{+}(x).\]
\end{theorem}

\begin{proof}[Proof of Theorem~\ref{thm:main-bridge-lower}]
If \(J_{r,m}^{+}(x)=\infty\), there is nothing to prove. Otherwise, fix
\(\varepsilon_0>0\). Choose \(\nu\in\mathcal P(\mathcal W_r)\), supported on
\(\mathcal G_r\), such that \(\int g_r\,d\nu>(r+1)x\) and
\((H(\nu\mid\pi_r)-\log p_b)/(r+1)\le J_{r,m}^{+}(x)+\varepsilon_0\).
Choose \(\alpha>0\) such that
\(\int g_r\,d\nu>(r+1)x+2\alpha\).

Since \(\mathcal W_r\) is finite and \(\nu\) is supported on \(\mathcal G_r\),
we may choose types \(\tau_k\in\mathcal P(\mathcal W_r)\), with denominator
\(k\), such that \(\tau_k\to\nu\), \(\tau_k(\mathcal G_r)=1\), and
\(\int g_r\,d\tau_k>(r+1)x+\alpha\) for all sufficiently large \(k\).

On the event \(E_k\cap\{L_k^{(r)}=\tau_k\}\), the seed event creates two
certified tokens in the top favorable state \(e_\star\). Moreover, all free
blocks belong to \(\mathcal G_r\). Hence repeated application of
Theorem~\ref{thm:bridge-lemma} to the forced word
\(W_k b^{(m)} W_{k-1} b^{(m)}\cdots b^{(m)}W_1\) gives
\[
N_{n_k}\ge 2\prod_{i=1}^k\Gamma(W_i),
\qquad
n_k^{-1}\log N_{n_k}
\ge n_k^{-1}\log 2+k n_k^{-1}\int g_r\,d\tau_k .
\]
Since \(n_k/k\to r+1\), the right-hand side is larger than \(x\) for all
sufficiently large \(k\). Thus
\(\mathbb P(n_k^{-1}\log N_{n_k}>x)\ge
\mathbb P(E_k\cap\{L_k^{(r)}=\tau_k\})\) for all sufficiently large \(k\).

By independence of the seed event, the bridge events, and the free blocks,
and by the standard method-of-types estimate on the finite alphabet
\(\mathcal W_r\),
\[
\mathbb P\bigl(E_k\cap\{L_k^{(r)}=\tau_k\}\bigr)
=
p_m p_b^{k-1}\,
\mathbb P\bigl(L_k^{(r)}=\tau_k\bigr),
\qquad
\lim_{k\to\infty}
k^{-1}\log \mathbb P\bigl(L_k^{(r)}=\tau_k\bigr)
=
-H(\nu\mid\pi_r).
\]
Using also \(n_k/k\to r+1\), the fixed seed cost \(\log p_m\) and the missing
single bridge factor are negligible after division by \(n_k\). Hence
\[
\liminf_{k\to\infty}
n_k^{-1}\log\mathbb P\bigl(n_k^{-1}\log N_{n_k}>x\bigr)
\ge
\frac{\log p_b-H(\nu\mid\pi_r)}{r+1}
\ge
-J_{r,m}^{+}(x)-\varepsilon_0 .
\]
Letting \(\varepsilon_0\downarrow0\) gives the claim.
\end{proof}

\end{document}
